\documentclass[10pt,a4paper]{article}
\usepackage{amsmath,amssymb,amsthm,graphicx}
\title{Near-total statistical-distance lower bound for sampling prefix ranks}
\author{Emanuele Viola}
\date{doi:10.4230/LIPIcs.CCC.2023.26; CC BY4.0}
\begin{document}
\maketitle
\noindent\textit{Editorial note.} Original author: Emanuele Viola. Official CCC2023 article26, 2023-07-10, https://drops.dagstuhl.de/entities/document/10.4230/LIPIcs.CCC.2023.26. Creative Commons Attribution4.0, https://creativecommons.org/licenses/by/4.0/.\par\medskip
\noindent\textit{Editorial note.} Faithful editable transcription of official CC BY4.0 main Theorem4 and complete Section7 proof/Lemma23, complete separator/fixed-set/comet/entropy support, original definitions and bibliography. Original author mathematical expressions retained, including original 1-2\textasciicircum{}-t0 wording. Long formulas mechanically split for standalone readability. No mathematical corrections or new proofs.\par\medskip
\noindent\textit{Allowed standard prerequisites (editorial).}\begin{itemize}
\item Markov inequality, union bound, entropy and statistical distance. Standard elementary probability; all specialized fixed-set machinery is fully included from author AppendixA.
\end{itemize}\medskip
\paragraph{The model (Definition 1).}
We say that $f:W^N\to\Sigma^m$ is a depth-$q$ forest with word size $w$ and output alphabet $\Sigma$ if $W=\{0,1\}^w$ and $f=(f_1,f_2,\ldots,f_m)$ where each $f_i:W^N\to\Sigma$ is a depth-$q$ decision tree where the variables are over $W$, each internal node has $|W|$ children, and the leaves are labeled with elements from $\Sigma$.
The main goal of this paper is to show that a target distribution $S$ over $\Sigma^m$ is hard to sample by a low-depth forest. We measure the distance between distributions $X$ and $Y$ over $D$ using statistical (a.k.a. total variation, $L_1$) distance
\[\Delta(X,Y):=\max_{T\subseteq D}|\mathbb P[X\in T]-\mathbb P[Y\in T]|.\]
So the lower-bound goal is to show $\Delta(f(U_{W^N}),S)$ is large for any low-depth forest $f$, where $U_{W^N}$ is the uniform distribution over $W^N$. In general for a set $H$ we write $U_H$ for the uniform distribution over $H$, and simply $U$ when $H$ is clear from the context.
\paragraph{Definition 13 (Almost pairwise independence).}
Jointly distributed random variables $X,Y$ are $\epsilon$-independent if $(X,Y)$ is $\epsilon$-close in statistical distance to $(X,Y')$ where $Y'$ has the same distribution of $Y$ and is independent from $X$.
The min-entropy $H_\infty(X)$ of a random variable $X$ is $\min_a\log_2(1/\mathbb P[X=a])$.
\paragraph{Notation.}
To avoid clutter in the more technical exposition of the results, we adopt the convention that for a set $S$ we also denote by $S$ the uniform distribution $U_S$ over $S$. The meaning will be clear from the context. For example, we shall simply write $\Delta(f(W^N),S)$ for $\Delta(f(U_{W^N}),U_S)$.

\paragraph{Definition 2.}
For $x\in\{0,1\}^m$ define $\operatorname{Rank}(x)$ as the string $y\in\{0,1,\ldots,m\}^m$ where $y_i:=\sum_{j\leq i}x_j$ is the rank of $i$.
\paragraph{Example 3.}
$\operatorname{Rank}(0,1,0,1)=(0,1,1,2)$.

\paragraph{Theorem 14 (Cell-probe sampling separator).}
There exists an integer $c\geq1$ such that the following holds:
\textbf{Hypothesis:} Let $f:W^N\to\Sigma^m$ be a depth-$q$ forest with word size $w$. Let $\alpha\leq1/c$. Let $t_0,t_1,\ldots$ be a sequence of integers with $t_i\geq t_{i-1}\cdot cqw/\alpha$ for every $i$. Let $S\subseteq\Sigma^m$ be a set and suppose that $\Delta(f(W^N),S)\leq1-2^{-t_0}$ where $2^{-t_0}\geq\sqrt{8/|S|}$.
\textbf{Conclusion:} There exists $k$, $1\leq k\leq O(q/\alpha)$, $D_k\subseteq W^N$, and $t_k$ indices $T\subseteq[m]$ such that:
\begin{enumerate}
\item[(0)] $H_\infty(f(D_k))\geq H_\infty(S)-t_{k-1}\cdot O(qw/\alpha)^2$;
\item[(1)] The support of $f(D_k)$ is contained in $S$;
\item[(2)] For every $i,j\in T$ the random variables $(f_i(D_k),f_j(D_k))$ are $O(\alpha)$-independent.
\end{enumerate}
For example, we can set $t_i=m/w^{a(q/\alpha)^{-bi}}$ which for suitable $a,b$ and $q,w\leq\log m$ satisfies the hypothesis.

\paragraph{Theorem 4.}
Let $f:W^N\to\{0,1,\ldots,m\}^m$ be a depth-$q$ forest with word size $w\geq\log m$. Then $\Delta(f(U_{W^N}),\operatorname{Rank}(U_{\{0,1\}^m}))\geq1-2\cdot2^{-m/w^{O(q)}}$.
Throughout this paper, the notation $O(.)$ and $\Omega(.)$ denotes absolute constants.

\paragraph{3. Proof of the separator Theorem 14.}
First we need to understand what it means to have slightly non-trivial statistical distance. Let $P$ be a distribution over $\Sigma^m$. One way in which $P$ can have statistical distance $\leq1-\epsilon$ from $S$ is if $P$ is distributed like $S$ with probability $\epsilon$, and it is say fixed with probability $1-\epsilon$. In this case, $P$ has actually very high entropy $(\log_2|S|)$ conditioned on an event of probability $\epsilon$. The next lemma shows that this in fact always happens.
\paragraph{Lemma 15.}
Let $P$ be a distribution over $\Sigma^m$ and let $S\subseteq\Sigma^m$. Suppose that $\Delta(P,S)\leq1-\epsilon$, where $\epsilon\geq\sqrt{8/|S|}$. Then there is a subset $S_0\subseteq S$ of probability $\mathbb P_P[S_0]=\Omega(\epsilon)$ such that the distribution $P$ conditioned on $P\in S_0$ has min-entropy $\geq H_\infty(S)-O(\log1/\epsilon)$.
\begin{proof}
We also write $P$ for the random variable distributed according to $P$. Collect all the elements of $S$ in increasing order of mass according to $P$ until right before collecting cumulative mass $\epsilon/2$. Note we don't collect all of $S$, for else $\mathbb P[P\in S]\leq\epsilon/2$ and $\Delta(P,S)\geq1-\epsilon/2$, contradicting the hypothesis.
Let $\beta$ be the mass of the next element of $S$. Let $S_0$ be the collected elements, $S_1$ the rest of $S$, and $T$ the complement of $S$. By definition, $\mathbb P[P\in S_0]<\epsilon/2$, and so $\mathbb P[P\in S_1\cup T]\geq1-\epsilon/2$. Also for every $x\in S_1$ we have $\mathbb P[P=x]\geq\beta$ and so $|S_1|\leq\beta^{-1}$. Combining these bounds with the assumption we have
\[1-\epsilon\geq\Delta(P,S)\geq\mathbb P[P\in S_1\cup T]-\frac{|S_1|}{|S|}\geq1-\epsilon/2-\frac{\beta^{-1}}{|S|}\]
and so $\beta\leq2/(\epsilon|S|)$.
Because we did not include in $S_0$ an element of mass $\beta$, and we only stop when we reach $\epsilon/2$, the mass of $S_0$ is $\geq\epsilon/2-\beta\geq\epsilon/2-2/(\epsilon|S|)$. If $\epsilon\geq\sqrt{8/|S|}$ this mass is at least $\epsilon/4$.
For any $x\in S_0$ using the above bound on $\beta$ we obtain
\[\mathbb P[P=x\mid P\in S_0]=\frac{\mathbb P[P=x]}{\mathbb P[P\in S_0]}\leq\frac{\beta}{\epsilon/4}\leq\frac8{\epsilon^2|S|},\]
as desired.
\end{proof}
The above lemma allows us to ``forget'' about $S$ and focus on $f$. We need to show that we can restrict the input to a large subset such that many output trees are nearly independent. This is the content of the following theorem. To avoid having to think about infinite sets, in the remainder of the proof we set the input to the sampler to $W^s$ for an integer $s$. This is without loss of generality, since obviously any forest of fixed depth can only access a finite number of input words.
We define the (entropy) loss of a subset $D'\subseteq D$ to be $\log_2(|D|/|D'|)$. So if $D'$ contains half the elements of $D$ the loss is one.
\paragraph{Theorem 16.}
There exists an integer $c\geq1$ such that the following holds:
\textbf{Hypothesis:} Let $f:W^s\to\Sigma^m$ be a depth-$q$ forest with word size $w$. Let $\alpha\leq1/c$. Let $t_0,t_1,\ldots$ be a sequence of integers with $t_i\geq t_{i-1}\cdot cqw/\alpha$ for every $i$. Let $D\subseteq W^s$ be a set with loss $\leq t_0$.
\textbf{Conclusion:} There exists $k$, $1\leq k\leq O(q/\alpha)$, $D_k\subseteq D$, and $t_k$ indices $T\subseteq[m]$ such that:
\begin{enumerate}
\item[(1)] The loss of $D_k\subseteq D$ is $\leq t_{k-1}\cdot(qw/\alpha)^2$;
\item[(2)] For every $i,j\in T$ the random variables $f_i(D_k),f_j(D_k)$ are $O(\alpha)$-independent.
\end{enumerate}
Let us first show how this gives the separator Theorem 14.
\begin{proof}[Proof of Theorem 14 from Theorem 16]
We apply Lemma 15 to $P=f(U)$. Given $S_0\subseteq S$ from the lemma, we let $D\subseteq W^s$ be the preimage of $S_0$ according to $f$. By the lemma, $|D|/|W|^s\geq\Omega(2^{-t_0})$, that is, the loss of $D\subseteq W^s$ is $t_0+O(1)$. Moreover, $H_\infty(f(D))\geq\log|S|-O(t_0)$.
We now apply Theorem 16 to this set $D$ and the sequence $t_0+O(1),t_1,t_2\ldots$. We can adjust the constant $c$ so that this satisfies the hypothesis. The theorem gives $D_k\subseteq D$ with loss $\leq t_{k-1}\cdot(qw/\alpha)^2$.
Observe that the support of $f(D_k)$ is contained in $S$, because the support of $f(D)$ is $S_0\subseteq S$ and $D_k\subseteq D$.
To verify the bound on $H_\infty(f(D_k))$, note that
\[\mathbb P[f(D)=x]\geq\mathbb P[f(D_k)=x]|D_k|/|D|.\]
Taking inverses and then logs we obtain
\[\log(1/\mathbb P[f(D)=x])\leq\log(1/\mathbb P[f(D_k)=x])+\log(|D|/|D_k|).\]
The left-hand side is at least $H_\infty(f(D))\geq\log|S|-O(t_0)$. While $\log(|D|/|D_k|)\leq t_{k-1}\cdot(qw/\alpha)^2$. Hence,
\[\log(1/\mathbb P[f(D_k)=x])\geq\log|S|-O(t_0)-t_{k-1}\cdot(qw/\alpha)^2,\]
for any $x$. The result follows.
\end{proof}
\paragraph{3.1. Proof of Theorem 16.}
The main technical lemma is the following one, which is like Theorem 16 but the requirement of independence is replaced by others easier to work with.
\paragraph{Lemma 17.}
Theorem 16 holds if we replace (2) with:
\begin{enumerate}
\item[(2')] for every $i,j\in T$: the probability over $D_k$ that $f_i(D_k)$ and $f_j(D_k)$ don't make all distinct probes is $\leq\alpha$, and
\item[(2'')] there exists a product distribution $R$ over words (that is, the words are independent) such that for every depth-$2q$ tree $g$, $g(D_k)$ and $g(R)$ are $\alpha$-close.
\end{enumerate}
\paragraph{Lemma 18.}
(2') and (2'') in Lemma 17 imply (2) in Theorem 16.
\begin{proof}
Let $X=D_k$. Think of $(f_i(X),f_j(X))$ as the output of the tree $g$ obtained by appending $f_j$ to the leaves of $f_i$. Note that $g$ makes $2q$ probes, possibly repeated. By (2''), there is a product distribution $R$ such that $g(X)$ and $g(R)$ are $\alpha$-close. Also, the probability that $g$ repeats a probe over $X$ is $\alpha$-close to the probability that it repeats it over $R$. Here we use that this probability can be written as the probability that a tree $d$ of depth $2q$ outputs 1, and that the output distributions of $d$ over $X$ and $R$ are $\alpha$-close. (Tree $d$ is obtained from $g$ by replacing any repeated probe along any path with a leaf 1, and any other leaf with 0.)
By this and (2') the probability that $g$ repeats a probe over $R$ is $\leq2\alpha$. Because $R$ is product, as long as probes are not repeated the output distribution does not change if we answer the first $q$ probes with $R^1$ and the next $q$ probes with $R^2$ where $R^1,R^2$ are iid copies of $R$. This shows that $(f_i(X),f_j(X))$ is $O(\alpha)$-close to $(f_i(R^1),f_j(R^2))$. Using again (2''), we can replace each $R^i$ with $X^i$, where $X^1,X^2$ are iid copies of $X$. This gives that $(f_i(X),f_j(X))$ is $O(\alpha)$-close to $(f_i(X^1),f_j(X^2))$. Adjusting constants concludes the proof.
\end{proof}
\paragraph{3.2. Proof of Lemma 17.}
We shall be concerned with inputs in various subsets $X\subseteq D$. If an input word is constant for every $x\in X$ then it needs not be probed but can be ``hardwired'' in the trees. We shall assume that the trees are always simplified accordingly. We denote by $G(x,X)$ the total number of probes made by all trees on input $x$, where the trees are simplified with respect to $X$.
We use the following fixed-set lemma from [21]. We say that distributions $X$ and $Y$ are $\alpha$-indistinguishable by depth-$q$ trees if for any such tree $t$, the statistical distance between $t(X)$ and $t(Y)$ is $\leq\alpha$.
\paragraph{Lemma 19 ([21], Lemma 3.14).}
Let $B\subseteq W^s$ be a subset with loss $\leq b$, where $W=\{0,1\}^w$. There exists $B_1\subseteq B$ and a product distribution $R$ such that $B_1$ and $R$ are $\alpha$-indistinguishable by depth-$2q$ decision trees. Moreover, the loss of $B_1\subseteq W^s$ is $\leq b\cdot O(wq/\alpha)$.
For completeness we include the proof in Appendix A.
We begin by applying this lemma to $D$ obtaining $D_1\subseteq D$ with loss $t_0\cdot O(wq/\alpha)$. This is the beginning of Iteration 1.
Our goal is to show that at the beginning of Iteration $k$ there exists a subset $D_k\subseteq D$ enjoying the following properties:
\begin{enumerate}
\item[(1)] [in Theorem 16] the loss of $D_k$ is $\leq t_{k-1}\cdot(qw/\alpha)^2$,
\item[(2'')] [in Lemma 17] there exists a product distribution $R$ over words such that for every depth-$2q$ tree $g$, $g(D_k)$ and $g(R)$ are $\alpha$-close.
\item[(3)] $\max_{x\in D_k}G(x,D_k)\leq m(q-\alpha(k-1)/4)$.
\end{enumerate}
Note all these hold at the beginning of Iteration 1.
In an iteration, collect as many trees as possible such that for any two of them, the probability over $D_k$ that they intersect probes is $\leq\alpha$. If you have $t_k$, then (2') in Lemma 17 holds as well, concluding the proof.
Otherwise, you have a collection of $t_k$ trees such that any other tree will intersect a probe with one of those $t$ with probability $\geq\alpha$. Now we are going to use this to proceed to the next iteration, i.e., increase the value of $k$ by 1. Because $G$ is non-negative, Property (3) above implies that there can be at most $O(q/\alpha)$ iterations, as desired.
Write $Y$ for the $\leq t_kq$ words probed by the $t_k$ trees in $D_k$. This is done according to a canonical order, and is a valid definition because the first probe of a tree is fixed, the second is fixed once the answer to the first is, and so on.
\paragraph{Support size.}
Let $D_{k,y}$ be the inputs in $D_k$ with $Y=y$. We have
\[\mathbb E_Y[|D_k|/|D_{k,Y}|]=\sum_y\mathbb P[Y=y]\frac{|D_k|}{|D_{k,y}|}=\sum_y1\leq|W|^{t_kq}.\]
By Markov's inequality $\mathbb P_Y[|D_k|/|D_{k,Y}|\geq M]\leq|W|^{t_kq}/M$. And so with probability $\geq1-|W|^{t_kq}/M$ over $Y$ we have $|D_{k,Y}|\geq|D_k|/M$.
\paragraph{Intersection.}
For a tree $f_i$ let $I_i(x)$ equal 1 if on input $x$ tree $f_i$ intersects probes with at least one of the $t_k$ trees collected, and equal 0 otherwise. Note that for every input $x$ and fixing $y$ we have
\[G(x,D_{k,y})\leq G(x,D_k)-\sum_{i\in[m]}I_i(x).\]
Because $\mathbb P_{x\in D_k}[I_i(x)=1]\geq\alpha$ for every $i$, we have
\[\mathbb E_{x\in D_k}\left[\sum_{i\in[m]}I_i(x)\right]\geq\alpha m.\]
Because the inner sum is $\leq m$, by Markov's inequality we have that with probability $\geq\alpha/2$ over the choice of $Y$
\[\mathbb E_{x\in D_{k,Y}}\left[\sum_{i\in[m]}I_i(x)\right]\geq\alpha m/2,\]
and
\[\mathbb E_{x\in D_{k,Y}}G(x,D_{k,Y})\leq\mathbb E_{x\in D_{k,Y}}G(x,D_k)-\alpha m/2.\]
\paragraph{Combining the arguments.}
Selecting $M=2|W|^{t_kq}/\alpha$ above, and by a union bound, there is a value $\bar y$ so that
\[\mathbb E_{x\in D_{k,\bar y}}[G(x,D_{k,\bar y})]\leq\mathbb E_{x\in D_{k,\bar y}}G(x,D_k)-\alpha m/2;\]
and at the same time $|D_{k,\bar y}|\geq|D_k|\cdot\alpha|W|^{-t_kq}/2$. That is, we increase the loss by
\[\leq\log(1/\alpha)+wt_kq+1.\]
Recall that the loss of $D_k$ is $t_{k-1}\cdot(qw/\alpha)^2$.
Note that $D_{k,\bar y}$ is still uniform over its support, since it is $D_k$ conditioned on a particular choice for $\leq t_kq$ words. Even though the words are chosen adaptively in $D_k$, once we condition on a particular value, their locations are fixed.
\paragraph{Reducing $G$ for every input.}
By Markov's inequality,
\begin{align*}
&\mathbb P_{x\in D_{k,\bar y}}\bigl[G(x,D_{k,\bar y})\geq(\mathbb E_{x\in D_{k,\bar y}}G(x,D_k)-\alpha m/2)(1+\alpha/(4q))\bigr]\\
&\hspace{4em}\leq\frac1{1+\alpha/(4q)}\leq1-\alpha/(8q).
\end{align*}
Hence, for $\geq\alpha/(8q)$ fraction of the inputs $x$ in $D_{k,\bar y}$ we have
\begin{align*}
G(x,D_{k,\bar y})&\leq(\mathbb E_{x\in D_{k,\bar y}}G(x,D_k)-\alpha m/2)(1+\alpha/(4q))\\
&\leq m(q-\alpha(k-1)/4)(1+\alpha/(4q))-\alpha m/2\\
&\leq m(q-\alpha k/4),
\end{align*}
using (3) in the second inequality. Let $D'_{k,\bar y}$ be the set of these inputs. The above gives the desired bound on $\max_{x\in D'_{k,\bar y}}G(x,D'_{k,\bar y})$, and note that the loss of $D'_{k,\bar y}\subseteq D_{k,\bar y}$ is $\leq\log(8q/\alpha)$.
\paragraph{Fixed-set lemma.}
Finally, we apply the fixed-set lemma to $D'_{k,\bar y}$ to obtain $D_{k+1}$; this gives (2''). This application multiplies the loss by $O(wq/\alpha)$, bringing the loss of $D_{k+1}\subseteq D$ to
\[O(wq/\alpha)\cdot O(t_{k-1}\cdot(qw/\alpha)^2+\log(1/\alpha)+wt_kq+1).\]
We need this loss to be at most $t_k\cdot(qw/\alpha)^2$. Dividing by $wq/\alpha$ we need to verify that
\[O(t_{k-1}\cdot(qw/\alpha)^2)+O(\log1/\alpha)+O(wt_kq)+O(1)\leq t_k\cdot qw/\alpha.\]
We claim that each term on the left-hand side is at most one-fourth of the right-hand side. For the first term we use the hypothesis that $t_k\geq t_{k-1}\cdot cqw/\alpha$ for a large enough $c$, and for the third we use that $\alpha\leq1/c$ and pick $c$ large enough. This gives (1).
Because $D_{k+1}\subseteq D'_{k,\bar y}$, the bound on $G$ still holds for $D_{k+1}$, and this gives (3).

\paragraph{A. Proof of the fixed-set Lemma 19.}
Begin with $R$ equal to the uniform distribution over $W^s$. If there are $q$ words and $q$ values such that the probability of getting those values in $B$ is larger than $(1+\alpha)/W^q$ then we fix them to those values, in both $B$ and $R$. Now we have subsets of $W^{s-q}$, the loss has decreased by an additive $\log_2 1/(1+\alpha)=\Omega(\alpha)$, and we repeat the process.
Because the initial loss was $b$, this process stops after $O(b/\alpha)$ iterations. In the end, the loss inside the final universe is at most $b$, since we never increase loss. With respect to the original universe, because we fixed $O(qb/\alpha)$ words, the loss is at most $O(wqb/\alpha)+b\leq b\cdot O(wq/\alpha)$.
Let $B_1$ and $R$ be the distributions when the process stops. Consider any tree $g:W^s\to\{0,1\}$ of depth $q$. Let $P$ be the collection of paths in $g$ leading to the output 1. Note that each path $p\in P$ corresponds to $q$ input words and $q$ values for them. Write $P_X[p]$ for the probability of following path $p$ under distribution $X$. By above we have $P_{B_1}[p]\leq(1+\alpha)/W^q=(1+\alpha)P_R[p]$.
Hence
\[\mathbb P[g(B_1)=1]=\sum_{p\in P}P_{B_1}[p]\leq\sum_{p\in P}(1+\alpha)P_R[p]=(1+\alpha)\mathbb P[g(R)=1].\]
And so in particular $\mathbb P[g(B_1)=1]\leq\mathbb P[g(R)=1]+\alpha$.
Repeating the argument with 0 and 1 swapped yields the lemma for trees with boolean alphabet. To prove the lemma for a tree $g'$ with arbitrary alphabet, reduce to the case of boolean alphabet in the following standard way. Suppose that the statistical distance between $g'(R)$ and $g'(B_1)$ is $>\alpha$. This means that there exists a set $T$ such that
\[|\mathbb P[g'(R)\in T]-\mathbb P[g'(B_1)\in T]|>\alpha.\]
Define tree $g$ with boolean output as $g(x):=1$ iff $g'(x)\in T$; note this just amounts to changing the labels of the leaves of $g'$. Now the left-hand side of the inequality above can be written as
\[|\mathbb P[g(R)=1]-\mathbb P[g(B_1)=1]|\]
and this contradicts the result for trees with boolean outputs and concludes the proof of the lemma.

\paragraph{4. Comets.}
In this section we define comets and prove a comet-finding lemma which will be used in our sampling lower bound.
\paragraph{Definition 20.}
A $d$-comet is a triple of indices $(i,j,k)$ from $[m]$ with $i<j<k$ such that $j-i\geq d(k-j)$. We call $(j,k)$ the head and $(i,j)$ the tail. The head length is $k-j$. A set of comets $\{(i_h,j_h,k_h)\}_h$ is disjoint if the intervals $[i_h,k_h]$ are disjoint.
\paragraph{Lemma 21 (Comet-finding).}
A subset of $\{1,2,\ldots,m\}$ of size $m/\ell^b$ contains $\geq m/\ell^{b+c+O(1)}$ disjoint $\ell^c$-comets where the head lengths are all in $[\ell^h,\ell^{h+1}]$ for some integer $h\leq b+c+O(1)$, for any $m$, $b\leq\ell$, $c\leq\ell$, and $\ell\geq\log m$.
\begin{proof}
Let $d=\ell^c$. First we claim that any subset of size $n:=d\log m+2$ contains a $d$-comet. Let the elements in the set be $a_1,a_2,\ldots$ in increasing order. If $(a_1,a_2,a_3)$ is not a $d$-comet then $a_3-a_2>(a_2-a_1)/d$, and so $a_3-a_1=a_3-a_2+a_2-a_1\geq(a_2-a_1)(1+1/d)$. Then again if $(a_1,a_3,a_4)$ is not a $d$-comet we have $a_4-a_3\geq(a_3-a_1)/d$ and so $a_4-a_1\geq(a_3-a_1)(1+1/d)\geq(a_2-a_1)(1+1/d)^2$. If we continue this way $n-2$ times, we obtain $a_n\geq(1+1/d)^{n-2}>m$, which is a contradiction.
Now divide the $t:=m/\ell^b$ elements of the given set into consecutive blocks of size $n$. By the previous paragraph, each block contains a comet. Hence we have $\geq t/n-1$ disjoint $d$-comets.
At least half of these comets have heads of length $\leq O(mn/t)=\ell^{b+c+O(1)}$, otherwise half the comets have heads longer than that, and we run out of space. Let $C_i$ be the subset of these comets whose head length is in $[\ell^i,\ell^{i+1})$. We only need to consider $i\leq b+c+O(1)$. Hence, there exists $i=h$ and
\[\Omega\left(\frac tn\right)\frac1{b+c+O(1)}\geq\frac m{\ell^{b+c+O(1)}}\]
disjoint comets with head lengths in $[\ell^h,\ell^{h+1})$, using that both $b$ and $c$ are $\leq\ell$.
\end{proof}
\paragraph{5. A lemma about entropy.}
In this section we quickly recall a basic result about entropy which will be used in our sampling lower bounds. The entropy $H$ of a random variable $X$ is defined as $H(X):=\sum_x\Pr[X=x]\cdot\lg(1/\Pr[X=x])$. The conditional entropy $H(X\mid Y):=\mathbb E_{y\in Y}H(X\mid Y=y)$ (cf. Chapter 2 in [13]).
\paragraph{Lemma 22.}
Let $Z=(Z_1,\ldots,Z_k)$ where $Z_i$ is supported over a set $S_i$, and let $\sum_i\log|S_i|=M$. Suppose $H(Z)\geq M-a$. There is a set $G\subseteq[k]$ of size $|G|\geq k-a/\epsilon$ such that for any $i\in G$ we have
\[H(Z_i\mid Z_1Z_2\ldots Z_{i-1})\geq\log|S_i|-\epsilon.\]
In particular, $Z_i$ is $4\sqrt\epsilon$ close to uniform over $S_i$.
\begin{proof}
By the chain rule for entropy ([13], Equation 2.21)
\[\sum_{i\leq k}(\log|S_i|-H(Z_i\mid Z_1Z_2\ldots Z_{i-1}))\leq a.\]
Applying Markov inequality to the non-negative random variable $\log|S_i|-H(Z_i\mid Z_1Z_2\ldots Z_{i-1})$ (for random $i\in[k]$), we have
\[\mathbb P_{i\in[k]}[\log|S_i|-H(Z_i\mid Z_1Z_2\ldots Z_{i-1})\geq\epsilon]\leq a/(k\cdot\epsilon),\]
yielding the desired $G$.
The ``in particular'' part holds because conditioning reduces entropy: $H(Z_i)\geq H(Z_i\mid Z_1Z_2\ldots Z_{i-1})$ ([13], Equations 2.60 and 2.92) and then applying Pinsker's inequality ([14], Chapter 3; Exercise 17).
\end{proof}

\paragraph{7. Proof of Theorem 4.}
We can assume that $q\leq\log m$, for else the statistical statistical bound is trivial and the theorem is true. We apply the separator Theorem 14 with $\alpha=1/1000$ and the sequence
\[t_i:=m/w^{c_0(q/\alpha)^{-c_1i}},\]
for constants $c_0,c_1$ to be set later. For large enough $c_1$ this satisfies the hypothesis of the theorem that $t_i\geq t_{i-1}\cdot cqw/\alpha$. The hypothesis that $1-2^{-t_0}\geq\sqrt{8/|S|}=\sqrt{8/2^m}$ holds as well since $|S|=2^m$.
Let $k$ and $D_k$ be as given by the theorem. Let
\[X:=\operatorname{Rank}^{-1}f(D_k).\]
Note that this is a valid definition because $f(D_k)$ is in the range of $\operatorname{Rank}$, and the latter is 1-1. The separator theorem guarantees that $H_\infty(f(D_k))\geq m-t'_{k-1}$, where
\[t'_{k-1}:=t_{k-1}\cdot O(q^2w^2/\alpha^4).\]
Hence also $H(X)\geq m-t'_{k-1}$.
\paragraph{Comets.}
We now apply the comet-finding Lemma 21 to the $t_k=m/w^{c_0(q/\alpha)^{-c_1k}}$ trees given by the separator. For $c=1$, the lemma gives a set of
\[t'_k:=m/w^{c_0(q/\alpha)^{-c_1k}+O(1)}\]
disjoint $w$-comets. We shall only use that they are 100-comets, and their head lengths will not be relevant now. We want to find a comet whose outputs are ``sufficiently random.''
Define $a:=t'_{k-1}$ and $b:=t'_k$.
Partition $X$ into $b$ consecutive blocks, where each block contains exactly one comet and intersects no others. Let $Z_1,Z_2,\ldots,Z_b$ be the blocks, and let $|Z_i|=s_i$ with $\sum_i s_i=m$. Applying Lemma 22 we find $\geq b-a/\epsilon$ blocks $i$ such that $H(Z_i\mid Z_1Z_2\ldots Z_{i-1})\geq s_i-\epsilon$. We set $\epsilon=1/w$ (a sufficiently small constant would be enough), and we verify that $b-a/\epsilon\geq1$, yielding at least one block $i^*$ such that
\[H(Z_{i^*}\mid Z_1Z_2\ldots Z_{i^*-1})\geq s_{i^*}-\epsilon.\tag{1}\]
The inequality $b-a/\epsilon\geq1$ is true because we only lost $w^{O(1)}$ factors, so it suffices to make the constant $c_1$ large enough in the definition of the sequence $t_i$. Formally,
\[\frac ba=\frac{t_k}{t_{k-1}}\cdot\frac1{O(q^2w^2/\alpha^4)\cdot w^{O(1)}}\geq\frac{w^{c_1}}{w^{O(1)}}>w^{100}.\]
The inequalities holds for $c_1$ large enough and using $q\leq\log m$, $w\geq\log m$, $\alpha=\Theta(1)$.
\paragraph{Breaking correlation in the random comet.}
Hence we now have a comet $(p,i,j)$ that is contained in an interval $Z_{i^*}$ such that:
\begin{enumerate}
\item Equation 1 holds, and
\item $f_i(D_k),f_j(D_k)$ are $\alpha$-independent.
\end{enumerate}
The next lemma directly contradicts this and concludes the proof.
\paragraph{Lemma 23.}
Let $X_1,X_2,\ldots,X_m$ be $0$--$1$ random variables, and $(p,i,j)$ a $c$-comet for a sufficiently large $c$. Let $\ell:=i-p$ and $d:=j-i$.
Suppose that
\[H(X_{p+1},X_{p+2},\ldots,X_j\mid X_1,X_2,\ldots,X_p)\geq\ell+d-1/c.\]
Then there exists an integer $t$ such that
\begin{align*}
\mathbb P_X[\operatorname{Rank}(X)_j\geq t+\ell/2+d/2+c^{1/3}\sqrt d]&\geq1/10,\text{ and}\\
\mathbb P_X[\operatorname{Rank}(X)_i<t+\ell/2]&\geq1/10,\text{ but}
\end{align*}
\begin{multline*}
\mathbb P_X[\operatorname{Rank}(X)_j\geq t+\ell/2+d/2+c^{1/3}\sqrt d\ \wedge\ \operatorname{Rank}(X)_i<t+\ell/2]\\
\leq1/1000\quad(\ll1/10\cdot1/10).
\end{multline*}
\begin{proof}
Let us start with the last inequality, because we can prove it without getting our hands on $t$. The probability is at most
\[\mathbb P_X\left[\sum_{k=i+1}^jX_k\geq d/2+c^{1/3}\sqrt d\right].\]
By Pinsker's inequality ([14], Chapter 3; Exercise 17) the distribution of $X_{i+1},X_{i+2},\ldots,X_j$ is $4/\sqrt c$ close to the uniform $U_1U_2\ldots U_d$. Hence the above probability is
\[\leq\Pr_U\left[\sum_{k=1}^dU_k\geq d/2+c^{1/3}\sqrt d\right]+4/\sqrt c\leq1/2000+4/\sqrt c\leq1/1000.\]
where the second inequality follows from Chebyshev's inequality for sufficiently large $c$.
We now verify the first two inequalities in the conclusion of the lemma. Let $Y:=X_1,X_2,\ldots,X_p$ stand for the prefix, and $Z:=X_{p+1},X_{p+2},\ldots,X_j$ for the $\ell+d$ high-entropy variables. Let
\[A:=\{y\in\{0,1\}^p:H(Z\mid Y=y)\geq\ell+d-2/c\}\]
be the set of prefix values conditioned on which $Z$ has high entropy. We claim that $\mathbb P[Y\in A]\geq1/2$. This is because, applying Markov Inequality to the non-negative random variable $\ell+d-H(Z\mid Y=y)$ (for $y$ chosen according to $Y$),
\begin{align*}
\mathbb P[Y\notin A]&=\mathbb P_{y\in Y}[\ell+d-H(Z\mid Y=y)>2/c]\\
&\leq\mathbb E_{y\in Y}[\ell+d-H(Z\mid Y=y)]/(2/c)\\
&=(\ell+d-H(Z\mid Y))/(2/c)\leq(1/c)/(2/c)=1/2.
\end{align*}
Note that for every $y\in A$ we have, by definition, that the $(\ell+d)$-bit random variable $(Z\mid Y=y)$ has entropy at least $\ell+d-2/c$, and so by Pinsker's inequality ([14], Chapter 3; Exercise 17) the random variable $(Z\mid Y=y)$ is $(\epsilon:=4\sqrt{2/c})$-close to uniform over $\{0,1\}^{\ell+d}$. Therefore, for any subset $S\subseteq A$, the random variable
\[(Z\mid Y\in S)\text{ is }\epsilon\text{-close to uniform over }\{0,1\}^{\ell+d}.\tag{2}\]
Now define $t$ to be the largest integer such that
\[\mathbb P[Y\in A\wedge\operatorname{Rank}(X)_p\geq t]\geq1/4.\tag{3}\]
Since by definition of $t$ we have $\mathbb P[Y\in A\wedge\operatorname{Rank}(X)_p\geq t+1]<1/4$, we also have
\[\mathbb P[Y\in A\wedge\operatorname{Rank}(X)_p\leq t]\geq1/2-1/4=1/4.\tag{4}\]
We obtain the desired conclusions as follows, denoting by $U_1,U_2,\ldots,$ uniform and independent $0$--$1$ random variables. The first probability in the conclusion of the lemma is at least
\[\mathbb P[\operatorname{Rank}(X)_j\geq t+(\ell+d)/2+\sqrt\ell/c^{1/6}]\]
because $\ell\geq c\cdot d$.
Writing $\operatorname{Rank}(X)_j$ as the sum of the first $p$ bits and the rest, the above probability is at least
\begin{multline*}
\mathbb P\left[\sum_{k\leq\ell+d}Z_k\geq(\ell+d)/2+\sqrt\ell/c^{1/6}\mid Y\in A\wedge\operatorname{Rank}(X)_p\geq t\right]\\
\cdot\mathbb P[Y\in A\wedge\operatorname{Rank}(X)_p\geq t].
\end{multline*}
The second factor is $\geq1/4$ by (3). Also by (2) in the first factor we can replace the $Z_k$ with uniform bits changing the probability by at most $\epsilon$. Hence the first factor is at least
\[\mathbb P\left[\sum_{k\leq\ell+d}U_k\geq(\ell+d)/2+\sqrt\ell/c^{1/6}\right]-\epsilon.\]
In turn the probability is
\[\geq1/2-\sqrt\ell/c^{1/6}\cdot\Theta(1/\sqrt\ell)\geq1/2-\Theta(1/c^{1/6})\]
using an estimate of the central binomial coefficient provided e.g. in [13], Lemma 17.5.1.
Overall, the first probability in the conclusion of the lemma is
\[(1/2-\Theta(1/c^{1/6})-\epsilon)(1/4)\geq1/10\]
for large enough $c$.
We now turn to the second probability in the conclusion of the lemma. Proceeding in a similar way, this probability is at least
\begin{align*}
&\mathbb P\left[\sum_{k\leq\ell}Z_k<\ell/2\mid Y\in A\wedge\sum_kY_k\leq t\right]\cdot\mathbb P\left[Y\in A\wedge\sum_kY_k\leq t\right]\\
&\geq\left(\mathbb P\left[\sum_{k\leq\ell}U_k<\ell/2\right]-\epsilon\right)\cdot(1/4)\\
&\geq(1/2-\epsilon)\cdot(1/4)\geq1/10
\end{align*}
for all sufficiently large $c$. Here the second inequality uses (2) and (4).
\end{proof}

\begin{thebibliography}{99}
\bibitem[1]{published-1} Miklós Ajtai. Σ11 -formulae on finite structures. Annals of Pure and Applied Logic, 24(1):1–48, 1983.
\bibitem[2]{published-2} Andris Ambainis, Leonard J. Schulman, Amnon Ta-Shma, Umesh V. Vazirani, and Avi Wigderson. The quantum communication complexity of sampling. SIAM J. on Computing, 32(6):1570–1585, 2003.
\bibitem[3]{published-3} Sanjeev Arora and Boaz Barak. Computational Complexity. Cambridge University Press, 2009. A modern approach.
\bibitem[4]{published-4} Paul Beame, Russell Impagliazzo, and Srikanth Srinivasan. Approximating ACˆ0 by small height decision trees and a deterministic algorithm for \#ACˆ0SAT. In Proceedings of the 27th Conference on Computational Complexity, CCC 2012, Porto, Portugal, June 26-29, 2012, pages 117–125. IEEE Computer Society, 2012. doi:10.1109/CCC.2012.40.
\bibitem[5]{published-5} Chris Beck, Russell Impagliazzo, and Shachar Lovett. Large deviation bounds for decision trees and sampling lower bounds for AC0-circuits. Electronic Colloquium on Computational Complexity (ECCC), 19:42, 2012.
\bibitem[6]{published-6} Chris Beck, Russell Impagliazzo, and Shachar Lovett. Large deviation bounds for decision trees and sampling lower bounds for AC0-circuits. In IEEE Symp. on Foundations of Computer Science (FOCS), pages 101–110, 2012.
\bibitem[7]{published-7} Avraham Ben-Aroya, Dean Doron, and Amnon Ta-Shma. Explicit two-source extractors for near-logarithmic min-entropy. Electronic Colloquium on Computational Complexity (ECCC), 23:88, 2016. URL: http://eccc.hpi-web.de/report/2016/088, arXiv:TR16-088.
\bibitem[8]{published-8} Itai Benjamini, Gil Cohen, and Igor Shinkar. Bi-lipschitz bijection between the boolean cube and the hamming ball. In IEEE Symp. on Foundations of Computer Science (FOCS), 2014.
\bibitem[9]{published-9} Eshan Chattopadhyay, Jesse Goodman, and David Zuckerman. The space complexity of sampling. Electron. Colloquium Comput. Complex., page 106, 2021. URL: https://eccc. weizmann.ac.il/report/2021/106, arXiv:TR21-106.
\bibitem[10]{published-10} Eshan Chattopadhyay and David Zuckerman. Explicit two-source extractors and resilient functions. In ACM Symp. on the Theory of Computing (STOC), pages 670–683, 2016.
\bibitem[11]{published-11} Gil Cohen. Making the most of advice: New correlation breakers and their applications. In IEEE Symp. on Foundations of Computer Science (FOCS), pages 188–196, 2016. doi: 10.1109/FOCS.2016.28.
\bibitem[12]{published-12} Gil Cohen and Leonard J. Schulman. Extractors for near logarithmic min-entropy. Electronic Colloquium on Computational Complexity (ECCC), 23:14, 2016.
\bibitem[13]{published-13} Thomas Cover and Joy Thomas. Elements of Information Theory (Wiley Series in Telecom- munications and Signal Processing). Wiley-Interscience, 2006.
\bibitem[14]{published-14} Imre Csiszar and Janos Korner. Information Theory: Coding Theorems for Discrete Memoryless Systems. Academic Press, Inc., 1982.
\bibitem[15]{published-15} Anindya De and Thomas Watson. Extractors and lower bounds for locally samplable sources. In Workshop on Randomization and Computation (RANDOM), 2011.
\bibitem[16]{published-16} P. Erdős and R. Rado. Intersection theorems for systems of sets. J. London Math. Soc., 35:85–90, 1960. doi:10.1112/jlms/s1-35.1.85.
\bibitem[17]{published-17} Merrick L. Furst, James B. Saxe, and Michael Sipser. Parity, circuits, and the polynomial-time hierarchy. Mathematical Systems Theory, 17(1):13–27, 1984.
\bibitem[18]{published-18} Anna Gál and Peter Bro Miltersen. The cell probe complexity of succinct data structures. Theoretical Computer Science, 379(3):405–417, 2007.
\bibitem[19]{published-19} Alexander Golynski. Cell probe lower bounds for succinct data structures. In 20th ACM-SIAM Symp. on Discrete Algorithms (SODA), pages 625–634, 2009.
\bibitem[20]{published-20} Mika Göös and Thomas Watson. A lower bound for sampling disjoint sets. ACM Trans. Comput. Theory, 12(3):20:1–20:13, 2020. doi:10.1145/3404858.
\bibitem[21]{published-21} Aryeh Grinberg, Ronen Shaltiel, and Emanuele Viola.               Indistinguishability by ad- aptive procedures with advice, and lower bounds on hardness amplification proofs. In IEEE Symp. on Foundations of Computer Science (FOCS), 2018.                    Available at https://www.ccs.neu.edu/home/viola/papers/adaptivemajority.pdf.
\bibitem[22]{published-22} Johan Håstad. Computational limitations of small-depth circuits. MIT Press, 1987.
\bibitem[23]{published-23} Johan Håstad. On the correlation of parity and small-depth circuits. SIAM J. on Computing, 43(5):1699–1708, 2014.
\bibitem[24]{published-24} Russell Impagliazzo, William Matthews, and Ramamohan Paturi. A satisfiability algorithm for AC0 . In ACM-SIAM Symp. on Discrete Algorithms (SODA), pages 961–972, 2012.
\bibitem[25]{published-25} Stasys Jukna. Boolean Function Complexity: Advances and Frontiers. Springer, 2012.
\bibitem[26]{published-26} Xin Li. Improved two-source extractors, and affine extractors for polylogarithmic entropy. In IEEE Symp. on Foundations of Computer Science (FOCS), 2016.
\bibitem[27]{published-27} Mingmou Liu and Huacheng Yu. Lower bound for succinct range minimum query. In Konstantin Makarychev, Yury Makarychev, Madhur Tulsiani, Gautam Kamath, and Julia Chuzhoy, editors, ACM Symp. on the Theory of Computing (STOC), pages 1402–1415. ACM, 2020. doi:10.1145/3357713.3384260.
\bibitem[28]{published-28} Shachar Lovett and Emanuele Viola. Bounded-depth circuits cannot sample good codes. Computational Complexity, 21(2):245–266, 2012.
\bibitem[29]{published-29} Peter Bro Miltersen. Cell probe complexity – A survey, 1999. Invited talk/paper at Advances in Data Structures (Pre-conference workshop of FSTTCS’99).
\bibitem[30]{published-30} Ryan O’Donnell. Analysis of Boolean Functions. Cambridge University Press, 2014.
\bibitem[31]{published-31} Mihai Pǎtraşcu. Succincter. In 49th IEEE Symp. on Foundations of Computer Science (FOCS). IEEE, 2008.
\bibitem[32]{published-32} Mihai Patrascu and Mikkel Thorup. Time-space trade-offs for predecessor search. In Jon M. Kleinberg, editor, ACM Symp. on the Theory of Computing (STOC), pages 232–240. ACM, 2006. doi:10.1145/1132516.1132551.
\bibitem[33]{published-33} Mihai Pǎtraşcu and Emanuele Viola. Cell-probe lower bounds for succinct partial sums. In 21th ACM-SIAM Symp. on Discrete Algorithms (SODA), pages 117–122, 2010.
\bibitem[34]{published-34} Alexander A. Razborov. Pseudorandom generators hard for k-DNF resolution and polynomial calculus resolution. Ann. of Math., 181(2):415–472, 2015. doi:10.4007/annals.2015.181.2.1.
\bibitem[35]{published-35} Nathan Segerlind, Sam Buss, and Russell Impagliazzo. A switching lemma for small restrictions and lower bounds for k-DNF resolution. SIAM J. on Computing, 33(5):1171–1200, 2004.
\bibitem[36]{published-36} Mikkel Thorup. Mihai patrascu: Obituary and open problems. Bulletin of the EATCS, 109:7–13, 2013. URL: http://albcom.lsi.upc.edu/ojs/index.php/beatcs/article/view/163/176.
\bibitem[37]{published-37} Emanuele Viola. The Complexity of Distributions, Fall 2018 talk at the Simons Institute. https://www.youtube.com/watch?v=O78b085HE3w.
\bibitem[38]{published-38} Emanuele Viola. Bit-probe lower bounds for succinct data structures. SIAM J. on Computing, 41(6):1593–1604, 2012.
\bibitem[39]{published-39} Emanuele Viola. The complexity of distributions. SIAM J. on Computing, 41(1):191–218, 2012.
\bibitem[40]{published-40} Emanuele Viola. Extractors for turing-machine sources. In Workshop on Randomization and Computation (RANDOM), 2012.
\bibitem[41]{published-41} Emanuele Viola. Extractors for circuit sources. SIAM J. on Computing, 43(2):355–972, 2014.
\bibitem[42]{published-42} Emanuele Viola. Lower bounds for data structures with space close to maximum imply circuit lower bounds. Theory of Computing, 15:1–9, 2019. URL: https://theoryofcomputing.org/ articles/v015a018/v015a018.pdf.
\bibitem[43]{published-43} Emanuele Viola. Sampling lower bounds: boolean average-case and permutations. SIAM J. on Computing, 49(1), 2020. Available at https://www.ccs.neu.edu/home/viola/papers/ sampling-lower-bounds.pdf.
\bibitem[44]{published-44} Emanuele Viola, 2022. https://emanueleviola.wordpress.com/2022/09/14/myth-creation -the-switching-lemma/.
\bibitem[45]{published-45} Andrew Yao. Separating the polynomial-time hierarchy by oracles. In 26th IEEE Symp. on Foundations of Computer Science (FOCS), pages 1–10, 1985.
\bibitem[46]{published-46} Andrew Chi-Chih Yao. Should tables be sorted? J. of the ACM, 28(3):615–628, 1981. doi:10.1145/322261.322274.
\bibitem[47]{published-47} Huacheng Yu. Optimal succinct rank data structure via approximate nonnegative tensor decomposition. In Moses Charikar and Edith Cohen, editors, ACM Symp. on the Theory of Computing (STOC), pages 955–966. ACM, 2019. doi:10.1145/3313276.3316352.
\end{thebibliography}
\end{document}
